\documentclass{article}
\usepackage{graphicx} % Required for inserting images
\usepackage{ctex}
\usepackage{tikz}
\usetikzlibrary{calc,angles,quotes}
\usepackage{amsmath}
\newcommand{\p}{\partial}
\newcommand{\e}{\epsilon}
\newcommand{\sd}{\nabla\cdot}%散度
\newcommand{\xd}{\nabla \times}%旋度
\newcommand{\dv}[2]{\frac{\mathrm d #1}{\mathrm d #2}}
\newcommand{\dvv}[2]{\frac{\mathrm d^2 #1}{\mathrm d #2^2}}
\newcommand{\pdv}[2]{\frac{\partial #1}{\partial #2}}
\newcommand{\pdvv}[2]{\frac{\partial^2 #1}{\partial #2^2}}
\title{电动力学第二章习题解答}
\author{wangjiahe060511 }
\date{August 2026}

\begin{document}

\maketitle
\section{二维静电势的求解}
直接求解第二问，不管什么解析函数的性质。\\
先写出通解，\begin{align}
    \Phi=(A+B\ln r)+\sum_l(A_lr^l+B_l\frac{1}{r^l})(Csinl\theta+Dcosl\theta)
\end{align}
带入边界条件，发现凑不出来！这是为什么？\\
问题出在分离变量上。在分离$\theta$的时候，有一个平凡解$A+B\theta$，但是这个解不满足周期性，给舍掉了。但是在这里，$\theta$转一圈必然穿过边界，所以不必舍去。故其解为\begin{align}
\boxed{
\Phi(r,\theta)=
\begin{cases}
\dfrac{V_0}{\alpha}\theta,
& 0<\theta<\alpha,\\[8pt]
\dfrac{V_0(2\pi-\theta)}{2\pi-\alpha},
& \alpha<\theta<2\pi.
\end{cases}
}
\end{align}
\section{欧氏空间中的广义反演变换}
先在二维情况下猜测圆心与半径。对称性，新圆心一定在原点与老圆心之间的连线上。通过反演$x_0-r_0，x_0+r_0$两点,可得圆心为$$\frac{R^2x_0}{x_0^2-r_0^2}$$,半径为$$\frac{R^2r_0}{x_0^2-r_0^2}$$然后来证明这一点。
\begin{figure}[htbp]
\centering
\begin{tikzpicture}[
    scale=1.15,
    line cap=round,
    line join=round,
    every node/.style={font=\small},
    original/.style={
        draw=blue!70!black,
        line width=1.1pt
    },
    inverse/.style={
        draw=orange!85!black,
        line width=1.1pt
    },
    auxiliary/.style={
        draw=gray!70,
        dashed,
        line width=0.7pt
    }
]

%------------------------------------------------
% 只用于确定图中各点的位置
% d：原圆心到反演中心的距离
% r_0：原圆半径
% R^2：反演幂
%------------------------------------------------
\pgfmathsetmacro{\d}{5}
\pgfmathsetmacro{\rzero}{2}
\pgfmathsetmacro{\Rtwo}{6}
\pgfmathsetmacro{\ang}{110}

% 原圆直径端点
\pgfmathsetmacro{\ax}{\d-\rzero}
\pgfmathsetmacro{\bx}{\d+\rzero}

% A、B 的反演点
\pgfmathsetmacro{\apx}{\Rtwo/\ax}
\pgfmathsetmacro{\bpx}{\Rtwo/\bx}

% 反演圆圆心和半径
\pgfmathsetmacro{\cpx}{(\apx+\bpx)/2}
\pgfmathsetmacro{\rprime}{(\apx-\bpx)/2}

% 原圆上的任意点 P
\pgfmathsetmacro{\px}{\d+\rzero*cos(\ang)}
\pgfmathsetmacro{\py}{\rzero*sin(\ang)}

% P 的反演点 P'
\pgfmathsetmacro{\pnormtwo}{\px*\px+\py*\py}
\pgfmathsetmacro{\ppx}{\Rtwo*\px/\pnormtwo}
\pgfmathsetmacro{\ppy}{\Rtwo*\py/\pnormtwo}

% 各点坐标
\coordinate (O)  at (0,0);
\coordinate (A)  at (\ax,0);
\coordinate (C)  at (\d,0);
\coordinate (B)  at (\bx,0);
\coordinate (P)  at (\px,\py);

\coordinate (Ap) at (\apx,0);
\coordinate (Bp) at (\bpx,0);
\coordinate (Cp) at (\cpx,0);
\coordinate (Pp) at (\ppx,\ppy);

% 顶部说明
\node[gray!70!black] at (3.5,3.05)
    {沿同一射线反演};

\node at (3.5,2.65)
    {$OP\!\cdot\!OP'
      =OA\!\cdot\!OA'
      =OB\!\cdot\!OB'
      =R^2$};

% 中心轴与反演射线
\draw[gray!45,line width=0.7pt]
    (-0.15,0)--(7.25,0);

\draw[auxiliary]
    (O)--(P);

% 原截圆与像截圆
\draw[original]
    (C) circle[radius=\rzero];

\draw[inverse]
    (Cp) circle[radius=\rprime];

% 原圆和像圆中的直角三角形
\draw[original]
    (A)--(P)--(B);

\draw[inverse]
    (Ap)--(Pp)--(Bp);

% 两个直角标记
\pic[
    draw=green!50!black,
    line width=0.8pt,
    angle radius=4.5mm
] {right angle=A--P--B};

\pic[
    draw=green!50!black,
    line width=0.8pt,
    angle radius=2.8mm
] {right angle=Ap--Pp--Bp};

% 各点
\fill (O) circle (1.5pt);

\foreach \X in {A,B,P}{
    \fill[blue!70!black] (\X) circle (1.5pt);
}

\foreach \X in {Ap,Bp,Pp}{
    \fill[orange!85!black] (\X) circle (1.5pt);
}

% 原圆心 C
\draw[blue!70!black,line width=0.8pt]
    ($(C)+(-0.09,0)$)--($(C)+(0.09,0)$)
    ($(C)+(0,-0.09)$)--($(C)+(0,0.09)$);

% 点的标注
\node[below left]  at (O)  {$O$};
\node[below]       at (Bp) {$B'$};
\node[below]       at (Ap) {$A'$};
\node[below]       at (A)  {$A$};
\node[below]       at (C)  {$C$};
\node[below]       at (B)  {$B$};
\node[above left]  at (Pp) {$P'$};
\node[above right] at (P)  {$P$};

% 直径说明
\node[orange!85!black] at (\cpx,-0.85)
    {像直径 $A'B'$};

\node[blue!70!black] at (\d,-2.30)
    {原直径 $AB$};

% 几何结论
\node[gray!70!black,align=center] at (3.5,-2.78)
{
    每个过 $OC$ 的截面都得到同一个像直径 $A'B'$；\\
    这些像圆绕 $OC$ 轴旋转后组成球面
};

\end{tikzpicture}

\caption{Generated by ChatGPT}
\end{figure}
图中，$\Delta OAP\sim \Delta OP'A'$,$\Delta OBP\sim \Delta OP'B'$,角度关系即可知$\angle B'P'A'$为直角，从而AB是新直径。
\section{两个球面的反演}
利用上方结论，假设a在b左，设b在x=d处，对两球进行反演令坐标相等，为
\begin{equation}
    \frac{R^2d}{d^2-b^2}=\frac{R^2(d-c)}{(d-c)^2-a^2}
\end{equation}
化简得到
\begin{equation}
    d^2c-(c^2-b^2-a^2)d-cb^2=0
\end{equation}
判别式恒正，必有解。
\section{缺角球壳上均匀电荷分布产生的电势}
直接积分，不可避免的会产生椭圆函数。太傻逼。用一个聪明一点的办法。\\
根据轴对称性，如果其在轴上的解是$\Phi=\sum_l c_l r^l$,那么在全空间上的解就一定是$$c_l r^l P_l(\cos\theta)$$,这是由球谐函数在$\theta=0$处的取值决定的。在轴线上积分就非常的爽快。
答案是：($z < R$)：$$\Phi(z) = \frac{\sigma_0 R}{2\varepsilon_0 z} \left[ (R+z) - \sqrt{R^2 + z^2 - 2Rz\cos\alpha} \right]$$\\($z > R$)：$$\Phi(z) = \frac{\sigma_0 R}{2\varepsilon_0 z} \left[ (z+R) - \sqrt{R^2 + z^2 - 2Rz\cos\alpha} \right]$$
泰勒展开比对系数即可。（结果挺简单，但是比对系数非常烦人）球心电场方向始终指向挖去点，于是直接求导得到$$\boldsymbol{E}(0) = \frac{Q\sin^2\alpha}{16\pi\varepsilon_0 R^2} \hat{\boldsymbol{z}}$$
\section{两个导体球问题的求解}
为了量纲一致，令A球电势为$V_0$
\subsection{}
对每个感应电荷，其在另一个球的感应电荷易得。所以为
\begin{align}
\boxed{
d_n^B=\frac{b^2}{c-d_n^A},d_{n+1}^A=\frac{a^2}{c-d_n^B}
}\\
\boxed{Q_n^B=\frac{-bQ_n^A}{c-d_n^A},Q_{n+1}^A=\frac{-aQ_n^B}{c-d_n^B}}
\end{align}
\\初始值好说。\[d_1^A = 0, \quad Q_1^A = 4\pi \varepsilon_0 a V_0\]
\subsection{}
利用$$Q_i=C_{ij}\Phi_j$$带入题目中$\Phi_A=V_0,\Phi_B=0$,既有
\begin{align}
    \boxed{C_{AA}=\frac{\sum q_A}{V_0},C_{BA}=\frac{\sum q_B}{V_0}}
\end{align}
\subsection{}
直接带入即可，有\begin{align}
    \boxed{d_{n+1}^A=\frac{a^2}{c-\frac{b^2}{c-d_n^A}},{d_{n+1}^B=\frac{b^2}{c-\frac{a^2}{c-d_n^B}}}}
\end{align}
\subsection{}
$d_\infty$好搞。从递推关系出发，把$d_n^A,d_n^B$全都换成$d_\infty$,就可以解得
\begin{align}
    \boxed{d_\infty =\frac{a^2+c^2-b^2\pm \sqrt{(a^2+c^2-b^2)^2-4a^2c^2}}{2c}}
\end{align}
考虑到$$0<d_\infty<a,\pm\text{取负号}$$。\\
按照题目说的带入$d^A_n=d^A_{\infty}+X_N$,两边同时除以$X_nX_{n+1}$,就是有
\begin{align}
    \frac{c^2-cd_\infty-b^2}{X_n}-c=\frac{cd_\infty-a^2}{X_{n+1}}
\end{align}
高中数学常见的递推式，最终解得\begin{align}
\boxed{
d_n^A=
\frac{
2a^2c\left[
1-
\left(
\frac{
c^2-a^2-b^2-\sqrt{(a^2+c^2-b^2)^2-4a^2c^2}
}{
c^2-a^2-b^2+\sqrt{(a^2+c^2-b^2)^2-4a^2c^2}
}
\right)^{n-1}
\right]
}{
a^2+c^2-b^2+\sqrt{(a^2+c^2-b^2)^2-4a^2c^2}
-
\left(
a^2+c^2-b^2-\sqrt{(a^2+c^2-b^2)^2-4a^2c^2}
\right)
\left(
\frac{
c^2-a^2-b^2-\sqrt{(a^2+c^2-b^2)^2-4a^2c^2}
}{
c^2-a^2-b^2+\sqrt{(a^2+c^2-b^2)^2-4a^2c^2}
}
\right)^{n-1}
}
}
\end{align}

\begin{align}
\boxed{
d_n^B=
\frac{
b^2+c^2-a^2-\sqrt{(a^2+c^2-b^2)^2-4a^2c^2}
-
\left(
b^2+c^2-a^2+\sqrt{(a^2+c^2-b^2)^2-4a^2c^2}
\right)
\left(
\frac{
c^2-a^2-b^2-\sqrt{(a^2+c^2-b^2)^2-4a^2c^2}
}{
c^2-a^2-b^2+\sqrt{(a^2+c^2-b^2)^2-4a^2c^2}
}
\right)^n
}{
2c\left[
1-
\left(
\frac{
c^2-a^2-b^2-\sqrt{(a^2+c^2-b^2)^2-4a^2c^2}
}{
c^2-a^2-b^2+\sqrt{(a^2+c^2-b^2)^2-4a^2c^2}
}
\right)^n
\right]
}
}
\end{align}
\subsection{}
\begin{align}
\text{巨大根式} \longrightarrow 2ab \sinh u, \\
\text{复杂公比} \longrightarrow \alpha^2, \\
d_\infty^A \longrightarrow a\xi, \\
d_\infty^B \longrightarrow b\eta.
\end{align}
\subsection{}5.4已经解出。
\section{单位立方体电容的近似计算}
\subsection{}
电势积分显然为
\begin{align}
\boxed{\Phi(x_0,y_0,z_0)=\int_{-1/2}^{1/2}\int_{-1/2}^{1/2}\frac{\mathrm{d}x\,\mathrm{d}y}{\sqrt{(x_0-x)^2+(y_0-y)^2+z_0^2}}}
\end{align}
对称性为x,y可交换，且对于x,y,z均为偶函数
\subsection{}
这一问，一个很自然的问题是各点电势不一样，没办法直接算出一个统一的电势然后算出电容。这里我们用$E=\frac{Q^2}{2C}=\frac{\sum Q\phi}{2}$来算一个平均值。\\
\begin{align}
    \sum Q\phi =\int_S \sigma \phi_{tot}dS
\end{align}
其中$\phi_{tot}$由六个面的电势加起来组成。\\
很显然，这个只能数值计算。$$\phi_{tot}=\Phi(x,y,0)+\Phi(x,y,1)+4\Phi(0.5,y,z-0.5)$$
积分得到9.24758148，从而
\begin{align}
\boxed{C=\frac{6}{K}(4\pi \e _0 a)=0.64881829(4\pi \e _0 a)}
\end{align}
\subsection{}
本问纯几何关系。用$\Phi$太他妈长了，\\
\boxed{
\begin{pmatrix}
V_1 \\
V_2 \\
V_3
\end{pmatrix}
= \frac{1}{4\pi\varepsilon_0 a}
\begin{pmatrix}
2.86598408 & 1.58094191 & 1.36561365 \\
1.59169235 & 1.76706216 & 1.38299866 \\
1.37517306 & 1.38876789 & 1.68489726
\end{pmatrix}
\begin{pmatrix}
Q_1 \\
Q_2 \\
Q_3
\end{pmatrix}
}
\subsection{}
首先来看第五问。为了保证$$V_1=V_2=V_3=V_0$$,解得
\begin{align}
    Q_1:Q_2:Q_3=0.12015963:0.72927808:1
\end{align}
这是总电量。由于块数不同，转化成面电荷密度就是
\begin{align}
    \sigma_1:\sigma_2:\sigma_3=0.48064:0.72928:1
\end{align}
挺符合直觉。中间最低越靠近边缘越高。
于是在这个情况下，算出总电荷，给出电容为
$$0.64599349(4\pi\e_0 a)$$效果反而变差，说明这个把电势搞成一样确实是一个很烂的近似
\subsection{}
已经完成
\subsection{}
最后考虑一种更合理的非等距$3\times 3$划分：在距离每条边棱$a/4$的位置作分割线。这样，每个面仍然被分成三类子面积：

\begin{itemize}
    \item 第一类为中心的正方形，大小为$a/2\times a/2$；
    \item 第二类为靠近边棱的长方形，大小为$a/2\times a/4$；
    \item 第三类为靠近顶角的正方形，大小为$a/4\times a/4$。
\end{itemize}

三类子面积的总面积分别为
\begin{align}
    S_1
    &=
    6\left(\frac{a}{2}\right)^2
    =
    \frac{3}{2}a^2,
    \\
    S_2
    &=
    24\left(\frac{a}{2}\right)
    \left(\frac{a}{4}\right)
    =
    3a^2,
    \\
    S_3
    &=
    24\left(\frac{a}{4}\right)^2
    =
    \frac{3}{2}a^2.
\end{align}

计算各类子面积中心受到三类电荷的电势，并将第二类长方形对应的积分上下限作相应修改，可以得到
\begin{align}
\boxed{
\begin{pmatrix}
V_1\\
V_2\\
V_3
\end{pmatrix}
=
\frac{1}{4\pi\varepsilon_0a}
\begin{pmatrix}
2.26980183 & 1.47385436 & 1.29458116\\
1.53736038 & 1.69886792 & 1.30812492\\
1.32031177 & 1.36373137 & 1.88351808
\end{pmatrix}
\begin{pmatrix}
Q_1\\
Q_2\\
Q_3
\end{pmatrix}.
}
\end{align}

令立方体表面三类代表点的电势相等，即
\begin{align}
    V_1=V_2=V_3=V_0,
\end{align}
解上述线性方程组可得
\begin{align}
\boxed{
\begin{pmatrix}
Q_1\\
Q_2\\
Q_3
\end{pmatrix}
=
4\pi\varepsilon_0aV_0
\begin{pmatrix}
0.10292711\\
0.32143220\\
0.22604347
\end{pmatrix}.
}
\end{align}

因此三类总电荷之比为
\begin{align}
    Q_1:Q_2:Q_3
    =
    0.455342:1.421993:1.
\end{align}

考虑到三类子面积的总面积不同，面电荷密度之比为
\begin{align}
\boxed{
\sigma_1:\sigma_2:\sigma_3
=
0.455342:0.710996:1.
}
\end{align}

该结果仍然符合电荷向边棱和顶角聚集的物理图像。立方体携带的总电荷为
\begin{align}
Q
&=
Q_1+Q_2+Q_3
\\
&=
0.65040277
\left(
4\pi\varepsilon_0aV_0
\right).
\end{align}

所以这种非等距$3\times3$划分给出的电容估计为
\begin{align}
\boxed{
C
=
\frac{Q}{V_0}
=
0.65040277
\left(
4\pi\varepsilon_0a
\right).
}
\end{align}

与等距$3\times3$划分得到的
\begin{align}
C_{\mathrm{uniform}}
=
0.64599349
\left(
4\pi\varepsilon_0a
\right)
\end{align}
相比，非等距划分的结果更接近高精度数值
\begin{align}
C_{\mathrm{ref}}
=
0.6606782
\left(
4\pi\varepsilon_0a
\right).
\end{align}
这是因为边棱和顶角附近的面电荷密度变化更快，在这些区域使用更小的子面积，能够更好地描述电荷聚集现象。
\end{document}
